\subsection{A smooth relative corner retaining the tensor lattice}
The special-fiber construction in \fid{F40} extends through the
actual collision. The smaller generator below suffices; no
classification of Cohen--Macaulay modules over a tacnode is used.

\begin{finding}[\fid{F44}: relative smoothness, the actual action,
and two different Frobenius unit conditions]\label{f:relativeauslander}
Let $R$ be a commutative domain, let $0\ne\Delta\in R$, and set
\[
 C=R[x]/(x(x-\Delta)),\qquad L_1=C/(x),\qquad
 L_2=C/(x-\Delta),\qquad
 \Gamma=\End_C(C\oplus L_1).
\]
Then $\Gamma$ is free of rank five over $R$ and has a projective
relative diagonal resolution of length two. It is smooth and
proper over $R$; if $R$ is a regular DVR, its absolute global
dimension is at most three.

The smooth proper algebra
$\Lambda=R\times R\times M_2(\Gamma)$ has an idempotent $f$ such that
\[
 f\Lambda f=A=R\times R\times M_2(C).
\]
For the actual rank-three family, with $\Delta=z_1-z_2$ and
$x=z_1-z$, this is the specified algebra of \fid{F19}. There is an
$R$-free perfect right $\Lambda$-module $\widetilde M$ of rank
$96$ such that, for $\mathcal P=f\Lambda$,
\[
 \End_\Lambda(\mathcal P)=A,\qquad
 \operatorname{Hom}_\Lambda(\mathcal P,\widetilde M)
       \simeq M
\]
as the actual rank-$64$ right $A$-module of \fid{F21}.
Moreover
\[
 \operatorname{Ext}^1_\Lambda(\widetilde M,\widetilde M)
       \simeq(R/(\Delta))^{64}.
\]
This retains the collision order in a coherent extension inside
the smooth relative category.

In characteristic $p>3$ there is an explicit ring Frobenius lift
on $\Gamma$ inducing ordinary Frobenius on $C$, with unit
derived pullback structures on the displayed perfect objects.
At a nonunit $\Delta$, those structures are not unit after
forgetting to $R$. More generally, no lift in the entire
fixed-vertex family considered below makes $\widetilde M$
both a unit object for the larger-algebra pullback and a unit
object for the underlying $R$-linearization.
This statement does not negate the independent $R$-unit
structure on the original corner module constructed in
\fid{F36}.
\end{finding}

\begin{proof}
\par\smallskip\noindent\textbf{The actual endomorphism algebra.}\quad
Let $e_1,e_2$ be the idempotents for $C,L_1$. Define
$a:L_1\to C$ by $a(1)=x-\Delta$ and let $b:C\to L_1$ be
evaluation at zero. Multiplication is composition, with the
right factor applied first. Put $c=ab$. Then
\begin{equation}\label{eq:relativearrows}
 a=e_1ae_2,\qquad b=e_2be_1,\qquad
 ba=-\Delta e_2,\qquad c=ab,\qquad c^2=-\Delta c.
\end{equation}
These formulas describe the endomorphism algebra itself.
Indeed the four matrix Hom spaces are $C,R,Rb,Ra$, because
\[
 \operatorname{ann}_C(x)=R(x-\Delta).
\]
To check the annihilator, write an element as $u+vx$.
Its product with $x$ is $(u+v\Delta)x$, which vanishes
exactly when $u=-v\Delta$. Thus $e_1,e_2,a,b,c$ is an
$R$-basis of $\Gamma$. The relation $ba=-\Delta e_2$
reduces every longer alternating path to this basis, and
\[
 e_1\Gamma e_1=R[c]/(c(c+\Delta))\simeq C,\qquad
 x=c+\Delta e_1.
\]
For the actual family choose $x=z_1-z$. Then $L_1$ is
the multiplicity-$16$ branch $z=z_1$ and $L_2$ is the
multiplicity-$4$ branch $z=z_2$. In all named matrix formulas
the original coordinate is therefore
\begin{equation}\label{eq:relativez}
 z=z_1-x=z_2-c.
\end{equation}
This specifies the original generators, as well as the abstract
algebra type. At $\Delta=0$ the algebra is the one in
\fid{F40}, with the stated sign change in the corner coordinate.
After inverting $\Delta$, the two branches split and
$\Gamma\simeq M_2(R)\times R$.

\par\smallskip\noindent\textbf{The complete relative diagonal resolution.}\quad
All tensor products in this paragraph are over $R$. Set
\[
 Q_0=(\Gamma e_1\otimes e_1\Gamma)
        \oplus(\Gamma e_2\otimes e_2\Gamma),
\]
\[
 Q_1=(\Gamma e_1\otimes e_2\Gamma)_a
        \oplus(\Gamma e_2\otimes e_1\Gamma)_b,\qquad
 Q_2=\Gamma e_2\otimes e_2\Gamma.
\]
These are finite projective modules over
$\Gamma\otimes_R\Gamma^{\mathrm{op}}$ and have $R$-ranks
$13,12,4$, respectively. Let $\mu:Q_0\to\Gamma$ be
multiplication and define
\[
 d_1(u\otimes v)_a=ua\otimes v-u\otimes av,\qquad
 d_1(u\otimes v)_b=ub\otimes v-u\otimes bv,
\]
\[
 d_2(u\otimes v)=(ub\otimes v)_a+(u\otimes av)_b.
\]
The subscripts identify the two arrow summands. These are
bimodule maps and $\mu d_1=0$. The scalar in the deformed
relation cancels:
\[
 d_1d_2(u\otimes v)=uba\otimes v-u\otimes bav
     =-\Delta u\otimes v+u\otimes\Delta v=0.
\]

We prove exactness over $R$ directly. The left ideal bases are
$\Gamma e_1=\langle e_1,b,c\rangle$ and
$\Gamma e_2=\langle e_2,a\rangle$; the right ideal bases are
$e_1\Gamma=\langle e_1,a,c\rangle$ and
$e_2\Gamma=\langle e_2,b\rangle$.
The image of $d_1$ contains
\[
 \begin{gathered}
 a\otimes e_2-e_1\otimes a,\qquad
 a\otimes b-e_1\otimes c,\qquad
 b\otimes e_1-e_2\otimes b,\qquad
 c\otimes e_1-a\otimes b,\\
 b\otimes a+\Delta e_2\otimes e_2,\qquad
 b\otimes c+\Delta e_2\otimes b,\qquad
 c\otimes a+\Delta a\otimes e_2,\qquad
 c\otimes c+\Delta a\otimes b.
 \end{gathered}
\]
Use the pivot coordinates
\[
 a\otimes e_2,\ e_1\otimes c,\ e_2\otimes b,\ c\otimes e_1,\
 b\otimes a,\ b\otimes c,\ c\otimes a,\ c\otimes c.
\]
Their $8$-by-$8$ minor is triangular in two blocks, with
diagonal entries $1,-1,-1,1,1,1,1,1$. Its determinant is
a unit, irrespective of $\Delta$. Hence their span $V$
is a free direct summand of $Q_0$, and $Q_0/V$ is free of
rank five. Multiplication annihilates $V$ and induces a
surjection $Q_0/V\to\Gamma$ between free modules of rank
five. A surjective square matrix over a commutative ring
has unit determinant, so this map is an isomorphism.
It follows that $V=\ker\mu=\operatorname{im}d_1$.

The $a$-component of $d_2$ uses $u\mapsto ub$, which sends
$e_2,a$ to the two basis vectors $b,c$. Coordinate projection
gives a left inverse. Thus $d_2$ is a split $R$-module
injection with free quotient of rank eight. The induced
surjection $Q_1/\operatorname{im}d_2\to V$ is an
isomorphism by the same square-matrix argument. We have
therefore proved the entire exact sequence
\begin{equation}\label{eq:relativediagonal}
 0\longrightarrow Q_2\xrightarrow{d_2}Q_1
   \xrightarrow{d_1}Q_0\xrightarrow{\mu}\Gamma
   \longrightarrow0.
\end{equation}
This proves relative smoothness. The finite free
$R$-module $\Gamma$ is proper over $R$.
The resolution proof works over any commutative base;
the domain and nonzerodivisor hypotheses enter the later
Ext and Frobenius arguments.

If $R$ is a regular DVR, tensor
\eqref{eq:relativediagonal} on the left with any right
$\Gamma$-module $N$. As a left-module resolution of the
projective module $\Gamma$, the bounded resolution is split
exact, so tensoring is exact. Its terms are finite sums of
$Ne_i\otimes_R e_j\Gamma$. Every $R$-module has projective
dimension at most one; since $e_j\Gamma$ is $R$-free and
projective over $\Gamma$, each such term has projective
dimension at most one over $\Gamma$. A length-two resolution
by these terms gives $\operatorname{pd}_\Gamma N\le3$.
No sharpness assertion for this bound is required.

\par\smallskip\noindent\textbf{The exact lift of the tensor lattice.}\quad
Put
\[
 \Lambda=R\times R\times M_2(\Gamma),\qquad
 f=(1,1,I_2\otimes e_1),\qquad \mathcal P=f\Lambda.
\]
Finite products and matrix Morita equivalence preserve
smoothness and properness, so $\Lambda$ is smooth proper
and $R$-free of rank $22$. Left multiplication gives
\[
 \End_\Lambda(\mathcal P)=f\Lambda f
       =R\times R\times M_2(C)=A.
\]
This identification uses \eqref{eq:relativez}.

Write $\mathsf P_i=e_i\Gamma$. Left multiplication by
$a$ sends the basis $e_2,b$ to the independent vectors
$a,c$. Hence
\begin{equation}\label{eq:relativeS}
 0\longrightarrow\mathsf P_2\xrightarrow{a}\mathsf P_1
      \longrightarrow S\longrightarrow0
\end{equation}
is exact. The quotient $S$ is $R$-free of rank one;
$e_1$ acts by one and $e_2,a,b,c$ by zero.
In particular $S$ is perfect of projective dimension at
most one. Its corner and that of $\mathsf P_2$ are
\[
 \mathsf P_2e_1=Rb\simeq L_1,\qquad Se_1=S\simeq L_2.
\]
Indeed $bx=bc+\Delta b=-\Delta b+\Delta b=0$,
whereas $x$ acts on $S$ by $\Delta$.

For a right $\Gamma$-module $N$, let
$\operatorname{Row}(N)=(N,N)$ be the standard right
$M_2(\Gamma)$-Morita module. Define
\begin{equation}\label{eq:relativelift}
 \widetilde M=\epsilon_+^{20}\oplus\epsilon_-^4
       \oplus\operatorname{Row}(\mathsf P_2)^{16}
       \oplus\operatorname{Row}(S)^4.
\end{equation}
It is $R$-free of rank $24+16\cdot4+4\cdot2=96$
and perfect over $\Lambda$, of projective dimension
at most one. Taking the corner gives exactly
\[
 \widetilde M f=
 \epsilon_+^{20}\oplus\epsilon_-^4\oplus N_1^{16}
       \oplus N_2^4=M.
\]
Evaluation at $f$ identifies
$\operatorname{Hom}_\Lambda(f\Lambda,\widetilde M)$
with $\widetilde M f$. Precomposition by left multiplication
in $f\Lambda f$ is the original right $A$-action.
Thus both the named algebra and the tensor lattice are
recovered over the family, including the collision.
At the closed fiber this rank-$96$ perfect lift differs
from the rank-$104$ projective lift chosen in \fid{F40};
both have the same actual rank-$64$ corner.

\par\smallskip\noindent\textbf{A coherent extension preserving the order of collision.}\quad
Applying $\operatorname{Hom}_\Gamma(-,\mathsf P_2)$ to
\eqref{eq:relativeS} gives
\[
 Rb\xrightarrow{-\Delta}Re_2,
\]
because composition is $ba=-\Delta e_2$.
As $\Delta$ is a nonzerodivisor,
\[
 \operatorname{Hom}_\Gamma(S,\mathsf P_2)=0,\qquad
 \operatorname{Ext}^1_\Gamma(S,\mathsf P_2)=R/(\Delta).
\]
The opposite derived Hom vanishes:
$\mathsf P_2$ is projective and
$\operatorname{Hom}_\Gamma(\mathsf P_2,S)=Se_2=0$.
Both self-Hom rings are $R$, and both positive self-Ext
groups vanish, using \eqref{eq:relativeS} for $S$.
Consequently exactly the $4\cdot16$ oriented cross
pairs in \eqref{eq:relativelift} contribute:
\[
 \operatorname{Ext}^1_\Lambda(\widetilde M,\widetilde M)
       =(R/(\Delta))^{64}.
\]
There is no extra factor from matrix Morita equivalence.
If $\Delta=ut^m$ over a DVR, with $u$ a unit, this
group has length $64m$, hence length $128$ for the actual
order-two collision. The original corner Ext in \fid{F25}
has both orientations and is $(R/(\Delta))^{128}$,
of length $256$ when $m=2$. These two Ext groups are
not identified.

\par\smallskip\noindent\textbf{The relative exceptional pair and the corner limitation.}\quad
The two self-Hom complexes of $S,\mathsf P_2$ are $R$,
the reverse derived Hom vanishes, and the triangle
\eqref{eq:relativeS} recovers $\mathsf P_1$. Thus they
generate. The functors from $\operatorname{Perf}(R)$
defined by these two objects have adjoints using derived
Hom and its $R$-dual. These complexes are perfect over
$R$: all objects involved are perfect over the
$R$-proper algebra $\Gamma$. Therefore
\[
 \operatorname{Perf}(\Gamma)
       =\langle\operatorname{Perf}(R),\operatorname{Perf}(R)\rangle.
\]
On each finite-field fiber the same known additivity and
Frobenius-enriched realization as in \fid{F40} apply:
\[
 U(\Gamma)_{\Q}=U(k)_{\Q}^{\oplus2},\qquad
 U(\Lambda)_{\Q}=U(k)_{\Q}^{\oplus4}.
\]
More explicitly, the invoked primary statements are
\cite[Example~12.3, p.~24; \S6.2, equation~(6.9), p.~16;
\S7, equations~(7.1)--(7.4), p.~17]{TabuadaWeil}.
The normalized TP spectrum on such a $\Lambda$-fiber
is four eigenvalues one in even degree and no odd part.
No fraction field of an arbitrary trait is being assumed
finite.

The corner functor
\[
 \operatorname{Perf}(C)\longrightarrow\operatorname{Perf}(\Gamma),
 \qquad X\longmapsto X\otimes_C^{\mathbf L}e_1\Gamma
\]
is fully faithful: its ambient derived right adjoint is
corner restriction, and the unit is the identity
$X\otimes_C^{\mathbf L}C=X$.
If $\Delta$ is a nonunit, this functor is not
right-admissible. Indeed the alternating resolution
with maps $x$ and $x-\Delta$ of $L_1$ gives
\[
 \operatorname{Tor}_{2j+1}^C(L_1,L_1)=R/(\Delta)
       \quad(j\ge0).
\]
The annihilator identities needed for this resolution
follow directly from the basis $1,x$. Thus $L_1$ is
not perfect. Since $\mathsf P_2e_1=L_1$, a right
adjoint on perfect categories would have to represent
$\operatorname{Hom}_{D(C)}(-,L_1)$ on perfect objects.
Testing $C[n]$ and its endomorphisms forces the
representing object to be $L_1$, a contradiction.
After inverting $\Delta$, the relation
$ba=-\Delta e_2$ makes $e_1$ full, and the corner
is instead a Morita equivalence.

\par\smallskip\noindent\textbf{A genuine noncommutative ring Frobenius lift.}\quad
Assume now $\operatorname{char}R=p>3$. Define the
$F_R$-semilinear ring map $\varphi:\Gamma\to\Gamma$ by
\begin{equation}\label{eq:relativefrobenius}
 \varphi(e_i)=e_i,\qquad
 \varphi(a)=a,\qquad
 \varphi(b)=\Delta^{p-1}b,\qquad F_R(r)=r^p.
\end{equation}
It respects the defining relation because
\[
 \varphi(b)\varphi(a)
       =-\Delta^p e_2=\varphi(-\Delta e_2).
\]
It is a ring map, not the operation taking every
noncommutative element to its $p$th power: for example
it fixes the nonzero square-zero arrow $a$.
On the commutative corner,
\[
 \varphi(c)=\Delta^{p-1}c=c^p,\qquad
 \varphi(x)=\Delta^{p-1}(c+\Delta e_1)=x^p.
\]
Here $c^2=-\Delta c$, $x^2=\Delta x$, and $p$ is odd.
The induced map is ordinary scheme Frobenius, including
on the actual coordinate $z=z_2-c$. Extend $\varphi$
to $\Lambda$ entrywise and on its scalar factors; it
fixes $f$.

For right modules write
$\mathbf L\varphi^*N=N\otimes_{\Gamma,\varphi}^{\mathbf L}\Gamma$.
The fixed vertex idempotents give
$\mathbf L\varphi^*\mathsf P_i\simeq\mathsf P_i$.
The differential $a$ in \eqref{eq:relativeS} remains
$a$ under $\varphi$, so
$\mathbf L\varphi^*S\simeq S$, with no additional
derived cohomology. Thus the displayed perfect objects
$\mathcal P,\widetilde M$ have unit structures for
pullback over their larger algebras.

\par\smallskip\noindent\textbf{The independent unit condition over the base.}\quad
The induced semilinear map on the $R$-basis $e_2,b$
of $\mathsf P_2$ sends these vectors to
$e_2,\Delta^{p-1}b$. Its $R$-linearization is
\begin{equation}\label{eq:relativebaseunit}
 \operatorname{diag}(1,\Delta^{p-1}),
\end{equation}
which is not invertible at a nonunit $\Delta$.
The induced $R$-linearization on $S$ is the identity.
There are $32$ copies of $b$ in the sixteen row copies
of $\mathsf P_2$ in \eqref{eq:relativelift}. Therefore
the determinant on the rank-$96$ lift is
$\Delta^{32(p-1)}$, up to a unit.
The recovered rank-$64$ corner has the same defect.
Changing the unit algebra-pullback isomorphism composes
it with an automorphism of the target, whose underlying
$R$-determinant is a unit; it cannot repair this defect.

We can classify this tradeoff for every
$F_R$-semilinear ring endomorphism fixing $e_1,e_2$.
The two arrow spaces have rank one, so its values must be
\[
 a\longmapsto ua,\qquad b\longmapsto vb
       \quad(u,v\in R).
\]
The relation and the nonzerodivisor $\Delta$ force
\begin{equation}\label{eq:relativetradeoff}
 uv=\Delta^{p-1}.
\end{equation}
Conversely this equality defines such a ring map and
gives ordinary Frobenius on $C$.
The pullback of \eqref{eq:relativeS} has differential
$ua$. Since $u\ne0$, it is injective, and its cokernel
as an $R$-module is
\[
 R\oplus(R/uR)^2.
\]
It is isomorphic to the free rank-one module $S$ only
if $u$ is a unit; a unit $u$ also suffices by rescaling
the resolution. The module $\mathsf P_2$ always has a
unit algebra-pullback structure, but its $R$-linearization
is $\operatorname{diag}(1,v)$, up to a unit change,
and is invertible exactly when $v$ is a unit.
For a nonunit $\Delta$, \eqref{eq:relativetradeoff}
prevents both conditions simultaneously.

For the full $\widetilde M$, a nonunit $u$ contributes
actual $R$-torsion from its $S$-summands, so it cannot
even be isomorphic to its algebra pullback. If $u$ is
a unit, $v$ is a nonunit and its underlying
$R$-linearization has a nonunit determinant. This
proves the assertion for the entire fixed-vertex
family, with the indicated scope.
\end{proof}

The conclusion concerns this chosen smooth lift and its
larger-algebra unit structures. The original corner module
still has the independent $R$-unit semilinear structure for
its intrinsic center Frobenius constructed in \fid{F36}.
The calculation above prevents that base-unit property
from being obtained through the unit pullback structures
of this particular lift; it does not prohibit the
\fid{F36} construction or every possible smooth model.
The ring map $\varphi$, scheme Frobenius on the commutative
corner, algebra-pullback unit structures, and cyclotomic
Frobenius on TP are distinct constructions. No equality
with cyclotomic Frobenius, Weil realization of all
specified morphisms, preservation of both original Ext
orientations, or positive nonstandard basis follows.

% Separate paragraph for Part III.
\paragraph{\fid{F44} in Grothendieck's language.}
The affine module/sheaf equivalence, precisely
Stacks Lemma~26.7.5, tag~01IB, translates the actual
corner lattice into coherent algebra on the affine
center. Perfect complexes mean bounded finite
projective complexes, as in Stacks Definition~15.76.1,
tag~0657; the modern dg smoothness criterion is
the perfect diagonal
\cite[\S5.1, pp.~13--14]{TabuadaWeil}.
The explicit relative diagonal and coherent derived
Hom prove the smooth repair and retain the order of
collision, while the nonadmissible corner prevents
an unwarranted motivic direct-summand assertion.
The Frobenius computation separates extension of
scalars over the endomorphism algebra from
linearization over the base. SGA~6, I, 4.7, supplies the
historical perfect-complex framework, not a
source for this noncommutative construction.
The unpublished archive leaves are framework
anchors only; none is claimed to contain this
relative algebra, its action comparison, or its
Frobenius obstruction.

The representation language is SGA~3, I, 4.7.2--3 and
cote~90, scan~17. Cote~16, scans~48--49, provides the
audited conditional semisimplicity/Weil vocabulary, not
a theorem about this relative noncommutative algebra.
